\documentclass[12pt,a4paper]{article}


\usepackage[utf8]{inputenc}
\usepackage[brazil]{babel}

\usepackage{amssymb}
\usepackage[T1]{fontenc}

\usepackage{graphicx}
\usepackage{enumerate}
\usepackage[final]{pdfpages}
\usepackage{pdflscape}
\usepackage{tikz}
\usetikzlibrary{arrows.meta}

\usepackage{cite}
\usepackage{bm}
%\usepackage{indentfirst}

\usepackage{amsmath}
\usepackage{amsthm,amsfonts}
\usepackage{latexsym}

\usepackage[left=2.5cm,right=3cm,top=3.5cm,bottom=3cm]{geometry}

\setcounter{tocdepth}{2}
%------------------------------------------------------------------------------
\begin{document}


%-------------------------------CAPA----------------------------------------
%------------------------------------------------------------------------------
\begin{titlepage}

\begin{center}

\large{\textbf{Universidade de Brasília - UnB Gama}}\\[1cm]



{\large Relatório de Física 1 Experimental}

\vspace{2.3cm}
{\large }
\vspace{2.6cm}



{\LARGE \bf  Experimento VI --- Conservação do\\ Momento e Colisões} \vspace{2cm}

Por
%
\vspace{2cm}

{\bf	Renann de Oliveira Gomes – 200043030 \\
		Ana Jullia B. M. Lopes – 232024554 \\
		Eduardo Henrique - 252015024 \\}



%
\vspace*{\fill}

%
{Brasília-DF, julho de 2026}


\end{center}
\end{titlepage}




\newpage

%\setcounter{chapter}{1}

\section{Objetivos}
Analisar e estudar colisões unidimensionais entre dois carrinhos sobre um trilho de ar, com as
forças resistivas reduzidas. Verificar experimentalmente a validade do princípio da conservação
do momento linear em colisões elásticas e completamente inelásticas, bem como da conservação da
energia cinética (no caso elástico) e a fração de energia dissipada (no caso inelástico).



\section{Introdução Teórica}

\subsection{Colisões elásticas}
Uma colisão elástica em um sistema isolado é aquela na qual há conservação simultânea do momento
linear e da energia cinética. As colisões entre dois carrinhos no trilho de ar são
aproximadamente elásticas quando reduzimos as forças resistivas --- como a força de atrito, que
praticamente desaparece quando os carrinhos flutuam sobre o colchão de ar gerado entre eles e o
trilho. Ao examinarmos a colisão entre um carrinho $A$ e um carrinho $B$ em uma dimensão, o
carrinho $B$ é deixado em repouso e o carrinho $A$ é colocado em movimento retilíneo uniforme
(MRU) para colidir com ele. Como o momento total do sistema deve ser conservado, o momento total
inicial é igual ao momento total final:
\begin{equation}
\label{eq:p}
m_A v_{Ai} + m_B v_{Bi} = m_A v_{Af} + m_B v_{Bf}.
\end{equation}

A energia cinética total inicial é igual à energia cinética total final:
\begin{equation}
\label{eq:k}
\frac{1}{2}m_A v_{Ai}^2 + \frac{1}{2}m_B v_{Bi}^2 =
\frac{1}{2}m_A v_{Af}^2 + \frac{1}{2}m_B v_{Bf}^2.
\end{equation}

Deixando o carrinho $B$ inicialmente em repouso, $v_{Bi}=0$, e resolvendo o sistema formado pelas
equações~\eqref{eq:p} e~\eqref{eq:k}, obtêm-se as velocidades finais em termos da velocidade
inicial do carrinho $A$:
\begin{equation}
\label{eq:vAf}
v_{Af} = \frac{m_A - m_B}{m_A + m_B}\, v_{Ai},
\end{equation}
\begin{equation}
\label{eq:vBf}
v_{Bf} = \frac{m_A}{m_B}\left(v_{Ai} - v_{Af}\right).
\end{equation}
As equações~\eqref{eq:vAf} e~\eqref{eq:vBf} serão utilizadas para analisar as duas leis de
conservação envolvidas: momento e energia.

\subsection{Colisões completamente inelásticas}
Nesse tipo de colisão a energia cinética do sistema não se conserva, sendo conservado apenas o
momento linear. Colocando o carrinho $A$ para colidir com o carrinho $B$ inicialmente em repouso,
$v_{Bi}=0$, após a colisão o carrinho $A$ se cola ao carrinho $B$ e os dois passam a se
movimentar acoplados com a mesma velocidade. A partir da equação~\eqref{eq:p}, as velocidades
finais serão $v_{Af}=v_{Bf}=v_{ABf}$, dada por
\begin{equation}
\label{eq:vABf}
v_{ABf} = \frac{m_A}{m_A + m_B}\, v_{Ai}.
\end{equation}

A razão entre a energia cinética final,
\begin{equation}
\label{eq:kf}
K_f = \frac{1}{2}\left(m_A + m_B\right)v_{ABf}^2,
\end{equation}
e a energia cinética inicial,
\begin{equation}
\label{eq:ki}
K_i = \frac{1}{2}m_A v_{Ai}^2,
\end{equation}
resulta em
\begin{equation}
\label{eq:razao}
\frac{K_f}{K_i} = \frac{m_A}{m_A + m_B}.
\end{equation}
Em uma colisão totalmente inelástica parte da energia cinética inicial é sempre perdida na
colisão.

Em ambas as partes, a velocidade de cada carrinho foi obtida a partir do intervalo de tempo que a
placa de largura $\ell = 10~\mathrm{cm} = 0{,}10~\mathrm{m}$, fixada sobre o carrinho, leva para
atravessar o sensor do cronômetro:
\begin{equation}
\label{eq:vel}
v = \frac{\ell}{\Delta t} = \frac{0{,}10~\mathrm{m}}{\Delta t}.
\end{equation}


\section{Parte Experimental e Discussão}

\subsection{Material utilizado}
\begin{itemize}
	\item kit trilho de ar;
	\item 02 carrinhos;
	\item 02 fixadores em U com elásticos (para os choques elásticos);
	\item 01 pino para carrinho com agulha e 01 pino com massa aderente (para os choques
	totalmente inelásticos);
	\item kit cronômetro com dois sensores ($S_1$ e $S_2$).
\end{itemize}

\subsection{1ª Parte --- Colisão elástica}
Os dois carrinhos foram posicionados sobre o trilho de ar com os fixadores em U com elásticos.
Sobre cada carrinho fixou-se uma placa retangular de $10~\mathrm{cm}$, que ao passar pelos
sensores permite ao cronômetro registrar os intervalos de tempo e, assim, calcular as
velocidades. Selecionou-se a função F3 do cronômetro: o sensor $S_1$ registra o intervalo de
tempo $\Delta t_1$ do carrinho $A$ e o sensor $S_2$ o intervalo $\Delta t_2$ do carrinho $B$. O
carrinho $A$ recebeu um impulso, entrando em MRU para colidir com o carrinho $B$, mantido em
repouso. O impulso é dado por um gancho (massa suspensa $m_g$) ligado ao carrinho $A$ por um fio
que passa por uma roldana na extremidade do trilho: a força peso do gancho traciona o fio e
acelera o carrinho $A$ até que este atinja a velocidade de MRU. A montagem e as forças atuantes
estão esquematizadas na Figura~\ref{fig:montagem}. As massas medidas dos carrinhos estão na
Tabela~\ref{tab:massas-el}.

\begin{figure}[!h]
\centering
\begin{tikzpicture}[>={Latex[length=2.4mm]}, thick, line join=round, scale=1.0]

	% ---------------- Trilho de ar ----------------
	\fill[gray!12] (0,0) rectangle (12,0.4);
	\draw[very thick] (0,0) rectangle (12,0.4);
	\draw[very thick] (1.2,0) -- (1.2,-0.35);
	\draw[very thick] (10.8,0) -- (10.8,-0.35);
	\node[anchor=north] at (3.2,-0.12) {\small Trilho de ar};

	% ---------------- Carrinho A ----------------
	\draw[very thick, fill=white] (0.7,0.4) rectangle (2.0,1.15);
	\node at (1.35,0.77) {$A$};
	\draw (0.95,1.15) -- (0.95,2.5);                 % haste (mastro)

	% ---------------- Carrinho B ----------------
	\draw[very thick, fill=white] (5.4,0.4) rectangle (6.6,1.05);
	\node at (6.0,0.72) {$B$};

	% ---------------- Fios e roldanas ----------------
	\draw (0.95,2.5) -- (11.85,2.5);                 % fio horizontal
	\draw[very thick, fill=white] (4.2,2.68) circle (0.2);    % sensor 1 (S1)
	\node[anchor=south] at (4.2,2.9) {\small Sensor 1 ($S_1$)};
	\draw[very thick, fill=white] (8.0,2.68) circle (0.2);    % sensor 2 (S2), após o carrinho B
	\node[anchor=south] at (8.0,2.9) {\small Sensor 2 ($S_2$)};
	\draw[very thick, fill=white] (11.9,2.45) circle (0.22);  % roldana da extremidade
	\node[anchor=south] at (11.9,2.72) {\small roldana};
	\draw (12.12,2.45) -- (12.12,0.55);              % fio descendo ao gancho
	\draw[very thick, fill=white]
	      (12.12,0.55) -- (11.82,0.2) -- (11.92,-0.25) -- (12.32,-0.25) -- (12.42,0.2) -- cycle;
	\node[anchor=west] at (12.5,0.15) {\small gancho ($m_g$)};

	% ---------------- Forças e velocidades ----------------
	% Tração no carrinho A (ao longo do fio, para a roldana)
	\draw[->, blue!55!black] (0.95,2.5) -- (2.35,2.5) node[above, pos=0.65] {$\vec{T}$};
	% Velocidade inicial de A
	\draw[->, red!70!black] (2.0,0.9) -- (3.3,0.9) node[above, pos=0.75] {$\vec{v}_{Ai}$};
	% Normal e peso de A
	\draw[->] (1.6,0.77) -- (1.6,1.75) node[right] {$\vec{N}_A$};
	\draw[->] (1.35,0.77) -- (1.35,-0.35) node[right, pos=0.95] {$\vec{P}_A$};
	% Normal e peso de B
	\draw[->] (6.0,0.72) -- (6.0,1.7) node[right] {$\vec{N}_B$};
	\draw[->] (6.0,0.72) -- (6.0,-0.4) node[right, pos=0.95] {$\vec{P}_B$};
	\node[anchor=west] at (6.75,0.6) {\small $v_{Bi}=0$};
	% Tração e peso do gancho
	\draw[->, blue!55!black] (12.12,0.15) -- (12.12,1.45) node[right, pos=0.9] {$\vec{T}$};
	\draw[->] (12.12,0.15) -- (12.12,-1.0) node[right, pos=0.95] {$\vec{P}_g = m_g\,\vec{g}$};

\end{tikzpicture}
\caption{Montagem do sistema para a colisão elástica no trilho de ar. O gancho de massa $m_g$
traciona o fio ($\vec{T}$) que, passando pela roldana, acelera o carrinho $A$ até o MRU com
velocidade $\vec{v}_{Ai}$, colidindo com o carrinho $B$ em repouso ($v_{Bi}=0$). Também estão
indicadas as forças normais ($\vec{N}_A$, $\vec{N}_B$), os pesos dos carrinhos ($\vec{P}_A$,
$\vec{P}_B$) e o peso do gancho ($\vec{P}_g = m_g\,\vec{g}$).}
\label{fig:montagem}
\end{figure}

\begin{table}[!h]
	\centering
	\begin{tabular}{|c|c|} \hline \hline
	massa do carrinho $A$ (kg) & massa do carrinho $B$ (kg) \\ \hline \hline
	0,2395 & 0,271 \\ \hline \hline
	\end{tabular}
	\caption{Massas dos carrinhos na colisão elástica.}
	\label{tab:massas-el}
\end{table}

Realizaram-se $5$ colisões, registrando-se os intervalos de tempo do carrinho $A$
($\Delta t_1$, antes da colisão) e do carrinho $B$ ($\Delta t_2$, após a colisão), conforme a
Tabela~\ref{tab:tempos-el}.

\begin{table}[!h]
	\centering
	\begin{tabular}{|c|c|c|} \hline \hline
	Colisão & $\Delta t_1$ (s) & $\Delta t_2$ (s) \\ \hline \hline
	1 & 0,202 & 0,216 \\ \hline
	2 & 0,205 & 0,219 \\ \hline
	3 & 0,202 & 0,214 \\ \hline
	4 & 0,201 & 0,217 \\ \hline
	5 & 0,202 & 0,205 \\ \hline
	Tempo médio & 0,202 & 0,214 \\ \hline \hline
	\end{tabular}
	\caption{Intervalos de tempo para as colisões elásticas.}
	\label{tab:tempos-el}
\end{table}

A partir dos tempos médios e da equação~\eqref{eq:vel}, a velocidade inicial do carrinho $A$ e a
velocidade final do carrinho $B$ (medida diretamente) são
\[
v_{Ai} = \frac{0{,}10}{0{,}202} = 0{,}494~\mathrm{m/s}, \qquad
v_{Bf} = \frac{0{,}10}{0{,}214} = 0{,}467~\mathrm{m/s}, \qquad v_{Bi} = 0.
\]

Pela equação~\eqref{eq:vAf}, a velocidade final do carrinho $A$ é
\[
v_{Af} = \frac{m_A - m_B}{m_A + m_B}\, v_{Ai}
       = \frac{0{,}2395 - 0{,}271}{0{,}2395 + 0{,}271}\times 0{,}494
       = -0{,}030~\mathrm{m/s},
\]
cujo sinal negativo indica que o carrinho $A$ recua após a colisão, coerente com o fato de
$m_B > m_A$. Pela equação~\eqref{eq:vBf}, a velocidade final teórica do carrinho $B$ é
\[
v_{Bf}^{\text{teo}} = \frac{m_A}{m_B}\left(v_{Ai} - v_{Af}\right)
       = \frac{0{,}2395}{0{,}271}\left(0{,}494 - (-0{,}030)\right) = 0{,}464~\mathrm{m/s}.
\]
Comparando com o valor medido $v_{Bf} = 0{,}467~\mathrm{m/s}$, o erro relativo é de apenas
$0{,}71\%$, mostrando boa concordância entre a previsão teórica e a medida experimental.

Os valores de momento linear e energia cinética estão resumidos na
Tabela~\ref{tab:result-el}. O momento inicial e final são
\[
P_i = m_A v_{Ai} = 0{,}118~\mathrm{kg\,m/s}, \qquad
P_f = m_A v_{Af} + m_B v_{Bf} = 0{,}119~\mathrm{kg\,m/s},
\]
e o erro relativo entre eles é
\[
E = \frac{|P_i - P_f|}{P_i} = 0{,}75\% < 5\%.
\]
As energias cinéticas inicial e final são
\[
K_i = \frac{1}{2}m_A v_{Ai}^2 = 0{,}029~\mathrm{J}, \qquad
K_f = \frac{1}{2}m_A v_{Af}^2 + \frac{1}{2}m_B v_{Bf}^2 = 0{,}030~\mathrm{J},
\]
com erro relativo
\[
E = \frac{|K_i - K_f|}{K_i} = 1{,}41\% < 5\%.
\]

\begin{table}[!h]
	\centering
	\begin{tabular}{|l|c|} \hline \hline
	Grandeza & Valor \\ \hline \hline
	Velocidade inicial $v_{Ai}$ (m/s)        & 0,494 \\ \hline
	Velocidade inicial $v_{Bi}$ (m/s)        & 0,000 \\ \hline
	Velocidade final $v_{Af}$ --- eq.~\eqref{eq:vAf} (m/s) & $-0{,}030$ \\ \hline
	Velocidade final $v_{Bf}$ --- medida (m/s)   & 0,467 \\ \hline
	Velocidade final $v_{Bf}$ --- teórica eq.~\eqref{eq:vBf} (m/s) & 0,464 \\ \hline
	Momento inicial $P_i$ (kg\,m/s)          & 0,118 \\ \hline
	Momento final $P_f$ (kg\,m/s)            & 0,119 \\ \hline
	Erro relativo do momento                 & 0,75\% \\ \hline
	Energia cinética inicial $K_i$ (J)       & 0,029 \\ \hline
	Energia cinética final $K_f$ (J)         & 0,030 \\ \hline
	Erro relativo da energia                 & 1,41\% \\ \hline \hline
	\end{tabular}
	\caption{Resultados da colisão elástica.}
	\label{tab:result-el}
\end{table}

\paragraph{Discussão.} Como os erros relativos do momento ($0{,}75\%$) e da energia cinética
($1{,}41\%$) são bem inferiores a $5\%$, pode-se afirmar que tanto o momento linear quanto a
energia cinética foram conservados na colisão, dentro das incertezas experimentais. Isso confirma
o comportamento aproximadamente elástico das colisões no trilho de ar quando as forças resistivas
são reduzidas. Os pequenos desvios observados decorrem do atrito residual com o colchão de ar, de
perdas nos elásticos dos fixadores em U e das flutuações nos tempos medidos pelos sensores.

\subsection{2ª Parte --- Colisão completamente inelástica}
Retirou-se o fixador em U com elástico, fixando-se o pino com agulha no carrinho $A$ e o pino com
massa aderente no carrinho $B$. Manteve-se a placa de $10~\mathrm{cm}$ sobre cada carrinho e a
função F3 no cronômetro, repetindo-se o procedimento: o carrinho $A$ recebe um impulso e entra em
MRU para colidir com o carrinho $B$ em repouso. Após a colisão, a agulha do pino do carrinho $A$
se prende à massa aderente do carrinho $B$, e os dois se movem colados. As massas medidas estão
na Tabela~\ref{tab:massas-in}.

\begin{table}[!h]
	\centering
	\begin{tabular}{|c|c|} \hline \hline
	massa do carrinho $A$ (kg) & massa do carrinho $B$ (kg) \\ \hline \hline
	0,230 & 0,221 \\ \hline \hline
	\end{tabular}
	\caption{Massas dos carrinhos na colisão completamente inelástica.}
	\label{tab:massas-in}
\end{table}

Nessas condições realizaram-se $5$ colisões, registrando-se o intervalo de tempo do carrinho $A$
antes da colisão ($\Delta t_1$) e o intervalo de tempo do conjunto acoplado após a colisão
($\Delta t_2$), conforme a Tabela~\ref{tab:tempos-in}.

\begin{table}[!h]
	\centering
	\begin{tabular}{|c|c|c|} \hline \hline
	Colisão & $\Delta t_1$ (s) & $\Delta t_2$ (s) \\ \hline \hline
	1 & 0,158 & 0,323 \\ \hline
	2 & 0,157 & 0,332 \\ \hline
	3 & 0,161 & 0,378 \\ \hline
	4 & 0,164 & 0,343 \\ \hline
	5 & 0,166 & 0,366 \\ \hline
	Tempo médio & 0,161 & 0,348 \\ \hline \hline
	\end{tabular}
	\caption{Intervalos de tempo para as colisões inelásticas.}
	\label{tab:tempos-in}
\end{table}

A partir dos tempos médios e da equação~\eqref{eq:vel}, a velocidade inicial do carrinho $A$ e a
velocidade final do conjunto acoplado (medida diretamente) são
\[
v_{Ai} = \frac{0{,}10}{0{,}161} = 0{,}620~\mathrm{m/s}, \qquad
v_{ABf} = \frac{0{,}10}{0{,}348} = 0{,}287~\mathrm{m/s}, \qquad v_{Bi} = 0.
\]

Pela equação~\eqref{eq:vABf}, a velocidade final teórica do conjunto é
\[
v_{ABf}^{\text{teo}} = \frac{m_A}{m_A + m_B}\, v_{Ai}
       = \frac{0{,}230}{0{,}230 + 0{,}221}\times 0{,}620 = 0{,}316~\mathrm{m/s}.
\]
O erro relativo entre a velocidade medida e a teórica é
\[
E = \frac{\left|v_{ABf} - v_{ABf}^{\text{teo}}\right|}{v_{ABf}} = 10{,}22\%.
\]

O momento inicial e final são
\[
P_i = m_A v_{Ai} = 0{,}143~\mathrm{kg\,m/s}, \qquad
P_f = (m_A + m_B)\, v_{ABf} = 0{,}129~\mathrm{kg\,m/s},
\]
com erro relativo
\[
E = \frac{|P_i - P_f|}{P_i} = 9{,}27\%.
\]

As energias cinéticas inicial e final são
\[
K_i = \frac{1}{2}m_A v_{Ai}^2 = 0{,}044~\mathrm{J}, \qquad
K_f = \frac{1}{2}(m_A + m_B)\, v_{ABf}^2 = 0{,}019~\mathrm{J},
\]
e a quantidade de energia perdida na colisão é
\[
Q = K_f - K_i = -0{,}026~\mathrm{J}.
\]

Por fim, verificando a validade da equação~\eqref{eq:razao}, a razão experimental entre as
energias cinéticas é
\[
\left(\frac{K_f}{K_i}\right)_{\text{exp}} = \frac{0{,}019}{0{,}044} = 0{,}420,
\]
enquanto o valor teórico previsto pela equação~\eqref{eq:razao} é
\[
\left(\frac{K_f}{K_i}\right)_{\text{teo}} = \frac{m_A}{m_A + m_B}
       = \frac{0{,}230}{0{,}451} = 0{,}510.
\]

\begin{table}[!h]
	\centering
	\begin{tabular}{|l|c|} \hline \hline
	Grandeza & Valor \\ \hline \hline
	Velocidade inicial $v_{Ai}$ (m/s)                     & 0,620 \\ \hline
	Velocidade inicial $v_{Bi}$ (m/s)                     & 0,000 \\ \hline
	Velocidade final $v_{ABf}$ --- medida (m/s)           & 0,287 \\ \hline
	Velocidade final $v_{ABf}$ --- teórica eq.~\eqref{eq:vABf} (m/s) & 0,316 \\ \hline
	Erro relativo da velocidade final                     & 10,22\% \\ \hline
	Momento inicial $P_i$ (kg\,m/s)                       & 0,143 \\ \hline
	Momento final $P_f$ (kg\,m/s)                         & 0,129 \\ \hline
	Erro relativo do momento                              & 9,27\% \\ \hline
	Energia cinética inicial $K_i$ (J)                    & 0,044 \\ \hline
	Energia cinética final $K_f$ (J)                      & 0,019 \\ \hline
	Energia perdida $Q = K_f - K_i$ (J)                   & $-0{,}026$ \\ \hline
	$K_f/K_i$ experimental                                & 0,420 \\ \hline
	$K_f/K_i$ teórica --- eq.~\eqref{eq:razao}            & 0,510 \\ \hline \hline
	\end{tabular}
	\caption{Resultados da colisão completamente inelástica.}
	\label{tab:result-in}
\end{table}

\paragraph{Discussão.} No choque totalmente inelástico o erro relativo do momento foi de
$9{,}27\%$, um pouco acima do limite de $5\%$. Ainda assim, o momento pode ser considerado
aproximadamente conservado: o valor superior ao caso elástico é esperado, pois no acoplamento dos
carrinhos há maior dissipação (deformação da massa aderente, cravamento da agulha e maior
dispersão nos tempos $\Delta t_2$ medidos, evidenciada pela maior variação entre as cinco
medidas). Já a energia cinética \emph{não} se conserva: houve perda de $Q = -0{,}026~\mathrm{J}$,
o equivalente a cerca de $58\%$ da energia cinética inicial, convertida em deformação, calor e som
no impacto. A razão experimental $K_f/K_i = 0{,}420$ é da mesma ordem do valor teórico
$0{,}510$ previsto pela equação~\eqref{eq:razao}, confirmando qualitativamente sua validade; a
diferença é compatível com as perdas por atrito e com a incerteza na medida de $\Delta t_2$.


\section{Conclusão}
O experimento permitiu analisar as duas leis de conservação em colisões unidimensionais no trilho
de ar. Na colisão elástica, tanto o momento linear (erro de $0{,}75\%$) quanto a energia cinética
(erro de $1{,}41\%$) foram conservados dentro das incertezas experimentais, confirmando o
caráter aproximadamente elástico do choque quando as forças resistivas são reduzidas; além disso,
a velocidade final do carrinho $B$ prevista pela teoria ($0{,}464~\mathrm{m/s}$) concordou com a
medida ($0{,}467~\mathrm{m/s}$) com erro de apenas $0{,}71\%$.

Na colisão completamente inelástica, o momento linear manteve-se aproximadamente conservado
(erro de $9{,}27\%$), enquanto a energia cinética não se conservou, com perda de
$0{,}026~\mathrm{J}$. A razão medida entre as energias cinéticas final e inicial
($K_f/K_i = 0{,}420$) mostrou-se compatível com o valor teórico da equação~\eqref{eq:razao}
($0{,}510$). Dessa forma, verificou-se experimentalmente que o momento linear se conserva em
ambos os tipos de colisão, ao passo que a energia cinética só se conserva na colisão elástica,
cumprindo os objetivos do experimento.

%\section{Referências Bibliográficas}
\begin{thebibliography}{}
%\bibliography{basename of .bib file}

\bibitem{Wytler_1} Wytler C. Santos, {\it Roteiro do Experimento VI --- Conservação do Momento e
Colisões}, disponibilizado no portal UnB Aprender3.

\bibitem{Wytler_2} Wytler C. Santos, {\it Medidas, algarismos e erros}, texto disponibilizado
no portal UnB Aprender3.

\bibitem {Young} H.D. Young \& R.A. Freedman, {\it Física I, Sears \& Zemansky}, Pearson,
Addison Wesley, $12^a$ Ed., (2009).




\end{thebibliography}


\end{document}
